% sec4_7.tex — Section 4.7: Application: Interrelated Factor Demands

\section{Application: Interrelated Factor Demands}
\label{sec:4.7}

The \textbf{translog cost function} is a generalization of the Cobb-Douglas (log-linear)
cost function introduced in Section~1.7. An attractive feature of the translog specification
is that the cost-minimizing input demand equations, if transformed into cost
share equations, are linear in the logs of output and factor prices, with the coefficients
inheriting (a subset of) the cost function parameters characterizing the technology.
Those technology parameters can be estimated from a system of equations
for cost shares. Early applications of this idea include Berndt and Wood (1975)
and Christensen and Greene (1976). This section reviews their basic methodology.
The methodology is also applicable in the analysis of consumer demand.

To better focus on the issue at hand and to follow the practice of those original
studies, we will proceed under the assumption of conditional homoskedasticity.
Also, we will assume that the regressors in the system of share equations, which are
log factor prices and log output, are predetermined. As we argued in Section~1.7,
this assumption is not unreasonable for the U.S.\ electric power industry before the
recent deregulation. Consequently, the appropriate multiple-equation estimation
technique is the multivariate regression.

\subsection*{The Translog Cost Function}

We have already made a partial transition from the log-linear to translog cost function:
in the empirical exercise of Chapter~1, we entertained the idea of adding the
square of log output to the log-linear cost function (see Model~4). If this idea is
extended to the quadratic and cross terms in the logs of all the arguments, the log-linear
cost function with three inputs (1.7.4) becomes the translog cost function:
\begin{multline}
\log(C) = \alpha_0 + \sum_{j=1}^{3} \alpha_j \log(p_j)
+ \frac{1}{2} \sum_{j=1}^{3} \sum_{k=1}^{3}
\gamma_{jk} \log(p_j) \log(p_k) \\
+ \alpha_Q \log(Q) + \frac{1}{2} \gamma_{QQ} (\log(Q))^2
+ \sum_{j=1}^{3} \gamma_{jQ} \log(p_j) \log(Q) + \varepsilon.
\label{eq:4.7.1}
\end{multline}
Here, and as in the rest of this section, the observation subscript ``$i$'' is dropped.
In this expression, the term
\[
\frac{1}{2} \sum_{j=1}^{3} \sum_{k=1}^{3}
\gamma_{jk} \log(p_j) \log(p_k)
\]
is a quadratic form representing the second-order effect of factor prices. We can
assume, without loss of generality, that the $3 \times 3$ matrix of quadratic form
coefficients, $(\gamma_{jk})$, is symmetric:\footnote{Let $\mathbf{x}$ be an
$n$-dimensional vector. The associated quadratic form is $\mathbf{x}'\mathbf{A}\mathbf{x}$.
Since $\mathbf{x}'\mathbf{A}\mathbf{x} = \mathbf{x}'\mathbf{A}'\mathbf{x}$, we have
$\mathbf{x}'\mathbf{A}\mathbf{x} = \mathbf{x}'[(\mathbf{A} + \mathbf{A}')/2]\mathbf{x}$.
So if $\mathbf{A}$ is not symmetric, it can be replaced by the symmetric matrix
$(\mathbf{A} + \mathbf{A}')/2$ without changing the value of the quadratic form.}
\begin{equation}
\gamma_{jk} = \gamma_{kj} \quad (j, k = 1, 2, 3).
\label{eq:4.7.2}
\end{equation}

In the case of log-linear cost function examined in Section~1.7, the degree of
returns to scale can be calculated as the reciprocal of the elasticity of costs with
respect to output. If this definition is applied to the translog cost function, we have
\begin{equation}
\text{returns to scale} = \frac{1}{\partial \log(C) / \partial \log(Q)}
= \frac{1}{\alpha_Q + \gamma_{QQ} \log(Q)
+ \sum_{j=1}^{3} \gamma_{jQ} \log(p_j)}.
\label{eq:4.7.3}
\end{equation}

\subsection*{Factor Shares}

The link between the cost function parameters and factor demands is given by
\textbf{Shephard's Lemma} from microeconomics. Let $x_j$ be the cost-minimizing demand
for factor input $j$ given factor prices $(p_1, p_2, p_3)$ and output $Q$.
So $\sum_{j=1}^{3} p_j x_j = C$. The lemma states that
\begin{equation}
\frac{\partial C}{\partial p_j} = x_j.
\label{eq:4.7.4}
\end{equation}
Noting that
\[
\frac{\partial \log(C)}{\partial \log(p_j)}
= \frac{p_j}{C} \cdot \frac{\partial C}{\partial p_j},
\]
the lemma can also be written as stating that the logarithmic partial derivative of
the cost function equals the factor share, namely,
\begin{equation}
\frac{\partial \log(C)}{\partial \log(p_j)} = \frac{p_j x_j}{C}.
\label{eq:4.7.5}
\end{equation}

For the case of translog cost function given in~\eqref{eq:4.7.1}, the log partial
derivative is easy to calculate:
\begin{equation}
\frac{\partial \log(C)}{\partial \log(p_j)}
= \alpha_j + \sum_{k=1}^{3} \gamma_{jk} \log(p_k) + \gamma_{jQ} \log(Q)
\quad (j = 1, 2, 3).
\label{eq:4.7.6}
\end{equation}
Combining this with~\eqref{eq:4.7.5} and defining the cost shares
$s_j \equiv p_j x_j / C$, we obtain the following system of cost share equations
associated with the translog cost function:
\begin{equation}
s_j = \alpha_j + \sum_{k=1}^{3} \gamma_{jk} \log(p_k) + \gamma_{jQ} \log(Q)
\quad (j = 1, 2, 3).
\label{eq:4.7.7}
\end{equation}
This system of share equations is subject to the \textit{cross-equation restrictions}
that the coefficient of $\log(p_k)$ in the $s_j$ equation equals the coefficient of
$\log(p_j)$ in the $s_k$ equation for $j \neq k$. In the rest of this section, these
cross-equation restrictions will be called the symmetry restrictions, although they
are \textit{not} a consequence of the symmetry assumption~\eqref{eq:4.7.2}.
Those cross-equation restrictions are a consequence of calculus: you should verify
that, if you did not assume symmetry, the coefficient of $\log(p_2)$ in the $s_1$
equation, for example, would be $(\gamma_{12} + \gamma_{21})/2$, which would
equal the $\log(p_1)$ coefficient in the $s_2$ equation.

\subsection*{Substitution Elasticities}

In the production function literature, a great deal of attention has been paid to
estimating the elasticities of substitution (more precisely, the Hicks-Allen partial
elasticities of substitution) between various inputs.\footnote{See Section~9.2 of
Berndt (1991) for a historical overview.}
As shown in Uzawa (1962), the substitution elasticity between $j$ and $k$,
denoted $\eta_{jk}$, is related to the cost function $C$ by the formula
\begin{equation}
\eta_{jk} = \frac{C \cdot \dfrac{\partial^2 C}{\partial p_j \partial p_k}}
{\dfrac{\partial C}{\partial p_j} \cdot \dfrac{\partial C}{\partial p_k}}.
\label{eq:4.7.8}
\end{equation}

For the translog cost function, some routine calculus and Shephard's lemma
show that
\begin{equation}
\eta_{jk} = \begin{cases}
\dfrac{\gamma_{jk} + s_j s_k}{s_j s_k} & \text{for } j \neq k, \\[10pt]
\dfrac{\gamma_{jj} + s_j^2 - s_j}{s_j^2} & \text{for } j = k,
\end{cases}
\label{eq:4.7.9}
\end{equation}
where $s_j$ is the cost share of input $j$. In the special case of the log-linear cost
function where $\gamma_{jk} = 0$ for all $j, k$, we have $\eta_{jk} = 1$ for
$j \neq k$, verifying that the substitution elasticities are unity between any two
factor inputs for the Cobb-Douglas technology.

\subsection*{Properties of Cost Functions}

As is well known in microeconomics (see, e.g., Varian, 1992, Chapter~5), cost
functions have the following properties:

\begin{enumerate}
\item \textit{(Homogeneity)}\quad
Homogeneous of degree one in factor prices.

\item \textit{(Monotonicity)}\quad
Nondecreasing in factor prices, which requires that
$\partial C / \partial p_j \geq 0$ for all $j$.

\item \textit{(Concavity)}\quad
Concave in factor prices, which requires that the Hessian (the
matrix of second-order derivatives) of a cost function $C$, say
\[
\mathbf{H} \equiv (\partial^2 C / \partial p_j \partial p_k),
\]
be negative semidefinite. By~\eqref{eq:4.7.8}, concavity is satisfied if and only if
the matrix of substitution elasticities, $(\eta_{jk})$, is negative
semidefinite.\footnote{Let $\mathbf{H}$ be the Hessian of the cost function $C$
and $\mathbf{V}$ be a diagonal matrix whose $j$-th diagonal is $\partial C / \partial p_j$.
Then the matrix of substitution elasticities, $(\eta_{jk})$, given in~\eqref{eq:4.7.8}
can be written compactly as: $(\eta_{jk}) = C \cdot \mathbf{V}^{-1} \mathbf{H}
\mathbf{V}^{-1}$. Since $\mathbf{V}$ is a diagonal matrix,
$\mathbf{V}^{-1} \mathbf{H} \mathbf{V}^{-1}$ is negative semidefinite if and only
if $\mathbf{H}$ is.}
\end{enumerate}

For the case of translog cost function, these properties take the following form.

\begin{enumerate}
\item \textit{Homogeneity:}\quad
It is easy to verify (see a review question) that, with the symmetry
restrictions imposed on $(\gamma_{jk})$, the homogeneity condition can be expressed as
\begin{equation}
\begin{cases}
\alpha_1 + \alpha_2 + \alpha_3 = 1, \\
\gamma_{11} + \gamma_{12} + \gamma_{13} = 0, \\
\gamma_{12} + \gamma_{22} + \gamma_{23} = 0, \\
\gamma_{13} + \gamma_{23} + \gamma_{33} = 0, \\
\gamma_{1Q} + \gamma_{2Q} + \gamma_{3Q} = 0.
\end{cases}
\label{eq:4.7.10}
\end{equation}
Homogeneity will be imposed in our estimation of the share equations.

\item \textit{Monotonicity:}\quad
Given Shephard's lemma, monotonicity is equivalent to requiring
that the right-hand sides of the share equations~\eqref{eq:4.7.7} be nonnegative for
any combinations of factor prices and output. For any values of coefficients,
we can always choose factor prices and output so that this condition is violated.
Accordingly, we cannot impose monotonicity on the share equations. It
is possible, however, to check whether monotonicity is satisfied for the ``relevant
range,'' that is, if it is satisfied for all the combinations of factor prices and
output found in the data at hand; we will do this for the estimated parameters
below.

\item \textit{Concavity:}\quad
Concavity requires that the matrix of substitution elasticities given
by~\eqref{eq:4.7.9} be negative semidefinite for any combinations of cost shares.
It is possible to show that a necessary and sufficient condition for concavity is that
the $3 \times 3$ matrix $(\gamma_{jk})$ be negative semidefinite.\footnote{Proof:
Let $\mathbf{s}$ here be a $3 \times 1$ vector whose $j$-th element is $s_j$ and
$\mathbf{S}$ here be a diagonal matrix whose $j$-th element is $s_j$.
Then the matrix of substitution elasticities given in~\eqref{eq:4.7.9} can be written
as $(\eta_{jk}) = \mathbf{S}^{-1}[(\gamma_{jk}) + \mathbf{s}\mathbf{s}' -
\mathbf{S}]\mathbf{S}^{-1}$. For some combination of cost shares,
$\mathbf{s}\mathbf{s}' - \mathbf{S}$ is zero (for example, set
$\mathbf{s} = (1, 0, 0)'$). So a necessary condition for the matrix $(\eta_{jk})$
to be negative semidefinite is that the matrix $(\gamma_{jk})$ is negative
semidefinite. This condition is also sufficient, because the matrix
$\mathbf{s}\mathbf{s}' - \mathbf{S}$ is negative semidefinite and the sum of two
negative semidefinite matrices is negative semidefinite.}
In particular, the diagonal elements, $(\gamma_{11}, \gamma_{22}, \gamma_{33})$,
have to be nonpositive. The condition that the matrix $(\gamma_{jk})$ be negative
semidefinite is a set of inequality conditions on its elements.
Estimation while imposing inequality conditions is feasible (see Jorgenson,
1986, Section~2.4.5, for the required technique), but it will not be pursued here.
\end{enumerate}

\subsection*{Stochastic Specifications}

The translog cost function given in~\eqref{eq:4.7.1} has an (additive) error term,
$\varepsilon$, while the share equations derived from it have none. On our data, the
share equations do not hold exactly for each firm, so we need to allow for errors.
The usual practice is to add a random disturbance term to each share equation.
There are two ways to justify such a stochastic specification. One is optimization
errors: firms make random errors in choosing their cost-minimizing input
combinations. But then actual total costs will not be as low as what is indicated
in the cost function, because there is no room for optimization errors in the
derivation of the cost function. The other justification for errors is to allow the
$\alpha_j$ in the share equations to be stochastic and differ across firms.
The intercept in the share equation for $j$ would then be the mean of $\alpha_j$,
and the error term would be the deviation of $\alpha_j$ from its mean. But this
also means that the $\alpha_j$ in the cost function has to be treated as random.

In either case, therefore, internal consistency would prevent us from adding
the cost function to the share equations with additive errors to form a system of
estimation equations. In the rest of this section, we will focus on a system that
consists entirely of share equations. This, however, means that we can estimate
only a subset of the cost function parameters. Two parameters, $\alpha_Q$ and
$\gamma_{QQ}$, which are absent in factor demands, cannot be estimated.
Since the degree of returns to scale depends on them (see~\eqref{eq:4.7.3}),
it cannot be estimated either.

\subsection*{The Nature of Restrictions}

With the additive errors appended, the system of share equations~\eqref{eq:4.7.7}
is written out in full as
\begin{equation}
\begin{cases}
s_1 = \alpha_1 + \gamma_{11} \log(p_1) + \gamma_{12} \log(p_2)
+ \gamma_{13} \log(p_3) + \gamma_{1Q} \log(Q) + \varepsilon_1, \\
s_2 = \alpha_2 + \gamma_{21} \log(p_1) + \gamma_{22} \log(p_2)
+ \gamma_{23} \log(p_3) + \gamma_{2Q} \log(Q) + \varepsilon_2, \\
s_3 = \alpha_3 + \gamma_{31} \log(p_1) + \gamma_{32} \log(p_2)
+ \gamma_{33} \log(p_3) + \gamma_{3Q} \log(Q) + \varepsilon_3.
\end{cases}
\label{eq:4.7.11}
\end{equation}

This system has fifteen coefficients or parameters. The cross-equation restrictions
(that the $3 \times 3$ matrix of log price coefficients, $(\gamma_{jk})$, be symmetric)
and the homogeneity restrictions~\eqref{eq:4.7.10} form a set of eight restrictions,
so that the fifteen parameters can be described by seven free parameters.
To understand this, it is useful to decompose the restrictions into the following three groups.
\begin{align}
\text{adding-up restrictions:} \quad &
\begin{cases}
\alpha_1 + \alpha_2 + \alpha_3 = 1, \\
\gamma_{11} + \gamma_{21} + \gamma_{31} = 0, \\
\gamma_{12} + \gamma_{22} + \gamma_{32} = 0, \\
\gamma_{13} + \gamma_{23} + \gamma_{33} = 0, \\
\gamma_{1Q} + \gamma_{2Q} + \gamma_{3Q} = 0.
\end{cases}
\label{eq:4.7.12} \\[6pt]
\text{homogeneity:} \quad &
\begin{cases}
\gamma_{11} + \gamma_{12} + \gamma_{13} = 0, \\
\gamma_{21} + \gamma_{22} + \gamma_{23} = 0, \\
\gamma_{31} + \gamma_{32} + \gamma_{33} = 0,
\end{cases}
\label{eq:4.7.13} \\[6pt]
\text{symmetry:} \quad &
\begin{cases}
\gamma_{12} = \gamma_{21}, \\
\gamma_{13} = \gamma_{31}, \\
\gamma_{23} = \gamma_{32}.
\end{cases}
\label{eq:4.7.14}
\end{align}

These are eleven equations, but only eight of them are restrictive: given the
adding-up restrictions, one of the three homogeneity restrictions is implied by
the two others, and given the adding-up and homogeneity restrictions, two of the
three symmetry restrictions are implied by the other (you should verify this).

\subsection*{Multivariate Regression Subject to Cross-Equation Restrictions}

It is straightforward to impose the homogeneity restrictions,~\eqref{eq:4.7.13},
because they are not cross-equation restrictions. We can eliminate three parameters,
say $(\gamma_{13}, \gamma_{23}, \gamma_{33})$, from the system to obtain
\begin{equation}
\begin{cases}
s_1 = \alpha_1 + \gamma_{11} \log(p_1/p_3)
+ \gamma_{12} \log(p_2/p_3) + \gamma_{1Q} \log(Q) + \varepsilon_1, \\
s_2 = \alpha_2 + \gamma_{21} \log(p_1/p_3)
+ \gamma_{22} \log(p_2/p_3) + \gamma_{2Q} \log(Q) + \varepsilon_2, \\
s_3 = \alpha_3 + \gamma_{31} \log(p_1/p_3)
+ \gamma_{32} \log(p_2/p_3) + \gamma_{3Q} \log(Q) + \varepsilon_3.
\end{cases}
\label{eq:4.7.15}
\end{equation}

The system (with or without the imposition of homogeneity) has the special
feature that the sum of the dependent variables, $(s_1, s_2, s_3)$, adds up to unity for
all observations. This feature and the adding-up restrictions imply that the sum of
the error terms, $(\varepsilon_1, \varepsilon_2, \varepsilon_3)$, is zero for all
observations, so the $3 \times 3$ error covariance matrix,
$\boldsymbol{\Sigma} \equiv \Var(\varepsilon_1, \varepsilon_2, \varepsilon_3)$,
is singular. Recall that, in the multivariate regression model (which is the special
case of the SUR model described in Proposition~4.6), the $\mathbf{S}$ matrix defined
in~\eqref{eq:4.1.11} can be written as a Kronecker product:
$\mathbf{S} = \boldsymbol{\Sigma} \otimes \E(\mathbf{x}_i \mathbf{x}_i')$.
Thus $\mathbf{S}$ becomes singular, in violation of Assumption~4.5. The common
practice to deal with this singularity problem is to drop one equation from the
system and estimate the system composed of the remaining equations (see below
for more on this issue). The coefficients in the equation that was dropped can be
calculated from the coefficient estimates for the included two equations by using
the adding-up restrictions. The adding-up restrictions as well as homogeneity are
thus incorporated.

As mentioned above, given the adding-up restrictions and homogeneity, there
is effectively only one symmetry restriction. For example, if $(\gamma_{13},
\gamma_{23}, \gamma_{33})$ are eliminated from the system to incorporate homogeneity,
as in~\eqref{eq:4.7.15}, and if the third equation is dropped to incorporate the
adding-up restrictions, then the unique symmetry restriction can be stated as
$\gamma_{12} = \gamma_{21}$. If this cross-equation restriction is imposed, the
system becomes
\begin{equation}
\begin{cases}
s_1 = \alpha_1 + \gamma_{11} \log(p_1/p_3)
+ \gamma_{12} \log(p_2/p_3) + \gamma_{1Q} \log(Q) + \varepsilon_1, \\
s_2 = \alpha_2 + \gamma_{12} \log(p_1/p_3)
+ \gamma_{22} \log(p_2/p_3) + \gamma_{2Q} \log(Q) + \varepsilon_2.
\end{cases}
\label{eq:4.7.16}
\end{equation}

Since the regressors are predetermined and since the equations have common
regressors, the system can be estimated by the multivariate regression subject to
the cross-equation restriction of symmetry. This constrained estimation, if it is cast
in the ``common coefficient'' format explored in Section~4.6, can be transformed
into an unconstrained estimation. That is, the system~\eqref{eq:4.7.16} can be
written in the ``common coefficient'' format of
$\mathbf{y}_i = \mathbf{Z}_i \boldsymbol{\delta} + \boldsymbol{\varepsilon}_i$
(see~\eqref{eq:4.6.1p} on page~292), if we set (while still dropping the ``$i$''
subscript from $s_1$, $s_2$, $\log(p_1/p_3)$, $\log(p_2/p_3)$, and $\log(Q)$)
\begin{equation}
\mathbf{y}_i = \begin{bmatrix} s_1 \\ s_2 \end{bmatrix}, \quad
\mathbf{Z}_i = \begin{bmatrix}
1 & 0 & \log(p_1/p_3) & \log(p_2/p_3) & 0 & \log(Q) & 0 \\
0 & 1 & \log(p_1/p_3) & 0 & \log(p_2/p_3) & 0 & \log(Q)
\end{bmatrix},
\quad
\boldsymbol{\delta}' = \begin{bmatrix}
\alpha_1 & \alpha_2 & \gamma_{11} & \gamma_{12} & \gamma_{22} &
\gamma_{1Q} & \gamma_{2Q}
\end{bmatrix}.
\label{eq:4.7.17}
\end{equation}
So the multivariate regression subject to symmetry is equivalent to the (unconstrained)
random-effects estimation.

This random-effects estimation of the two-equation system produces consistent
estimates of the following seven parameters:
\[
\alpha_1, \alpha_2, \gamma_{11}, \gamma_{12}, \gamma_{22},
\gamma_{1Q}, \gamma_{2Q}.
\]
The rest of the fifteen share equation parameters,
\[
\alpha_3, \gamma_{13}, \gamma_{21}, \gamma_{23}, \gamma_{31}, \gamma_{32},
\gamma_{33}, \gamma_{3Q},
\]
can be calculated using the adding-up restrictions~\eqref{eq:4.7.12},
homogeneity~\eqref{eq:4.7.13}, and symmetry~\eqref{eq:4.7.14}.
For example, $\gamma_{33}$ can be calculated as
\[
\gamma_{33} = -\gamma_{31} - \gamma_{32}
= -\gamma_{13} - \gamma_{23}
= (\gamma_{11} + \gamma_{12}) + (\gamma_{12} + \gamma_{22})
= \gamma_{11} + 2\gamma_{12} + \gamma_{22}.
\]
So the point estimate of $\gamma_{33}$ is given by
\[
\hat{\gamma}_{33} = \hat{\gamma}_{11} + 2\hat{\gamma}_{12} + \hat{\gamma}_{22},
\]
where $(\hat{\gamma}_{11}, \hat{\gamma}_{12}, \hat{\gamma}_{22})$ are the
random-effects estimates. The standard error of $\hat{\gamma}_{33}$
can be calculated by applying the ``delta method'' (see Lemma~2.5) to the
consistent estimate of
$\avar(\hat{\gamma}_{11}, \hat{\gamma}_{12}, \hat{\gamma}_{22})$ obtained from the
random-effects estimation.

\subsection*{Which Equation to Delete?}

Does it matter which equation is to be dropped from the system? It clearly would
not if there were no cross-equation restrictions, because the multivariate regression
is numerically equivalent to the equation-by-equation OLS. To see if the numerical
invariance holds under the cross-equation restriction, let
$\boldsymbol{\Sigma}^*$ be the $2 \times 2$ matrix of error covariances for the
two-equation system obtained from dropping one equation from the three-equation
system~\eqref{eq:4.7.15}. It is the appropriate submatrix of the $3 \times 3$ error
covariance matrix $\boldsymbol{\Sigma}$. (For example, if the third equation is dropped,
then $\boldsymbol{\Sigma}^*$ is the leading $2 \times 2$ submatrix of the $3 \times 3$
matrix $\boldsymbol{\Sigma}$.) To implement the random-effects estimation, we need a
consistent estimate, $\widehat{\boldsymbol{\Sigma}}^*$, of $\boldsymbol{\Sigma}^*$.
Two ways for obtaining $\widehat{\boldsymbol{\Sigma}}^*$ are:

\begin{itemize}
\item (equation-by-equation OLS)\quad Estimate the two equations separately, thus
ignoring the cross-equation restriction, and then use the residuals to calculate
$\widehat{\boldsymbol{\Sigma}}^*$. Equivalently, estimate the \textit{three}
equations separately by OLS, use the three-equation residuals to calculate
$\widehat{\boldsymbol{\Sigma}}$ (an estimate of the $3 \times 3$ error covariance
matrix $\boldsymbol{\Sigma}$), and then extract the appropriate submatrix from
$\widehat{\boldsymbol{\Sigma}}$. If $\widehat{\boldsymbol{\Sigma}}^*$ is thus
obtained, then (as you will verify in the empirical exercise) the numerical invariance
is guaranteed; it does not matter which equation to drop.

\item Obtain $\widehat{\boldsymbol{\Sigma}}^*$ from some technique that exploits
the cross-equation restriction (such as the pooled OLS applied to the common
coefficient format). The numerical invariance in this case is not guaranteed.
\end{itemize}

The numerical invariance holds under equation-by-equation OLS, because the
$3 \times 3$ matrix $\widehat{\boldsymbol{\Sigma}}$, from which
$\widehat{\boldsymbol{\Sigma}}^*$ is derived, is singular in finite samples.
In large samples, as long as $\widehat{\boldsymbol{\Sigma}}^*$ is consistent,
and even if $\widehat{\boldsymbol{\Sigma}}^*$ is not from equation-by-equation OLS,
this relationship between $\widehat{\boldsymbol{\Sigma}}^*$ and
$\widehat{\boldsymbol{\Sigma}}$ holds asymptotically because $\boldsymbol{\Sigma}$
is singular. So the invariance holds asymptotically if
$\widehat{\boldsymbol{\Sigma}}^*$ is consistent. In particular, the asymptotic
variance of parameter estimates does not depend on the choice of equation to drop.

\subsection*{Results}

Nerlove's study examined in Chapter~1 was based on 1955 cross-section data for
U.S.\ electric utility companies. Christensen and Greene (1976) updated Nerlove's
data to 1970 for 99 firms, using data sources that allowed them to construct better
measures of the wage rate and fuel prices (see their Section~V for more details).
Their data set includes information on factor shares, so the share equations can be
estimated. The means and standard deviations for some of the variables in the data
are reported in Table~4.3, which shows that fuel accounts for a very large share of
the cost of electricity generation.

\begin{table}[htbp]
\centering
\caption{Simple Statistics (Sample Size $= 99$)}
\label{tab:4.3}
\medskip
\begin{tabular}{@{}lcccc@{}}
\toprule
& Output in & Labor & Capital & Fuel \\
& kilowatt hours & share & share & share \\
\midrule
Mean & 9.0 & 0.141 & 0.227 & 0.631 \\
Std.\ deviation & 10.3 & 0.059 & 0.062 & 0.095 \\
\bottomrule
\end{tabular}
\end{table}

We use the equation-by-equation OLS residuals to calculate
$\widehat{\boldsymbol{\Sigma}}^*$, so that the random-effects estimates are
numerically invariant to the choice of equation to drop. The estimates are reported
in Table~4.4, along with $\widehat{\boldsymbol{\Sigma}}$ derived from the
equation-by-equation OLS. The factor input index $j$ is 1 for labor, 2 for capital,
and 3 for fuel. Some (actually all) of the diagonal elements of the matrix of
estimated price coefficients, $(\hat{\gamma}_{jk})$, are positive.
So the matrix is not negative semidefinite, in violation of concavity.\footnote{This
``test'' of concavity can be formalized. Using the delta method, we can calculate
the asymptotic distribution of the characteristic roots of $(\hat{\gamma}_{jk})$.
The null hypothesis of concavity is that those characteristic roots are all nonpositive.}

Even if $(\hat{\gamma}_{jk})$ is not negative semidefinite, the associated substitution
elasticities given in the formula~\eqref{eq:4.7.9} may be negative semidefinite.
Since the formula is derived assuming no optimization errors, fitted cost shares,
rather than actual cost shares, should be used for the $s_j$ and $s_k$ in the formula.
The elasticities averaged over firms are shown in Table~4.5. (They are the
off-diagonal elements of the symmetric $3 \times 3$ substitution elasticity matrix
thus calculated.) The three factor inputs are substitutes because the substitution
elasticities are positive, but the degree of substitutability is far less than what is
assumed in the log-linear technology. Concavity is violated even in the ``relevant
range'': for 20 firms out of the 99 in the sample, the matrix of substitution
elasticities is not negative semidefinite. Monotonicity, in contrast, is satisfied in the
relevant range: fitted cost shares are nonnegative for all inputs for all firms in the sample.

\begin{table}[htbp]
\centering
\caption{Random-Effects Estimates}
\label{tab:4.4}
\medskip
\begin{tabular}{@{}lccc@{}}
\toprule
Parameter & Point estimate & Standard error & $t$-value \\
\midrule
$\alpha_1$ & $-0.132$ & $0.106$ & $-1.25$ \\
$\alpha_2$ & $0.318$ & $0.085$ & $3.75$ \\
$\alpha_3$ & $0.813$ & $0.094$ & $8.69$ \\
$\gamma_{11}$ & $0.084$ & $0.020$ & $4.19$ \\
$\gamma_{12}$ & $-0.023$ & $0.016$ & $-1.46$ \\
$\gamma_{13}$ & $-0.060$ & $0.015$ & $-3.92$ \\
$\gamma_{22}$ & $0.122$ & $0.020$ & $6.19$ \\
$\gamma_{23}$ & $-0.099$ & $0.017$ & $-5.75$ \\
$\gamma_{33}$ & $0.159$ & $0.023$ & $6.90$ \\
$\gamma_{1Q}$ & $-0.0211$ & $0.0025$ & $-8.55$ \\
$\gamma_{2Q}$ & $-0.0086$ & $0.0030$ & $-2.87$ \\
$\gamma_{3Q}$ & $0.0297$ & $0.0037$ & $7.98$ \\
\bottomrule
\end{tabular}
\bigskip

$\widehat{\boldsymbol{\Sigma}}$ by pooled OLS $= \begin{bmatrix}
0.00173 & -0.000171 & -0.00156 \\
-0.000171 & 0.00253 & -0.00236 \\
-0.00156 & -0.00236 & 0.00391
\end{bmatrix}$
\end{table}

\begin{table}[htbp]
\centering
\caption{Substitution Elasticities}
\label{tab:4.5}
\medskip
\begin{tabular}{@{}ccc@{}}
\toprule
Labor-Capital & Capital-Fuel & Labor-Fuel \\
\midrule
0.17 & 0.27 & 0.29 \\
\bottomrule
\end{tabular}
\end{table}
